\documentclass{article}
\usepackage[utf8]{inputenc}
\usepackage{amsmath}

\begin{document}

\section*{Notatki z wykladu -- granice ciagow}

Mowimy, ze ciag $(a_n)$ ma granice $g$, gdy dla kazdego $\varepsilon > 0$
istnieje takie $N$, ze dla wszystkich $n > N$ zachodzi
$|a_n - g| < \varepsilon$. Piszemy wtedy $\lim_{n \to \infty} a_n = g$.

\subsection*{Twierdzenie o trzech ciagach}

Jezeli $a_n \leq b_n \leq c_n$ dla prawie wszystkich $n$ oraz
$\lim a_n = \lim c_n = g$, to rowniez $\lim b_n = g$.

\subsection*{Przyklady}

\begin{enumerate}
  \item $\lim_{n \to \infty} \frac{1}{n} = 0$
  \item $\lim_{n \to \infty} \left(1 + \frac{1}{n}\right)^n = e$
  \item $\lim_{n \to \infty} \sqrt[n]{n} = 1$
\end{enumerate}

\subsection*{Cwiczenia}

Pokaz, ze ciag $a_n = \frac{n^2 + 3n}{2n^2 - 1}$ jest zbiezny i wyznacz
jego granice. Wskazowka: podziel licznik i mianownik przez $n^2$.

\end{document}
